\section{Introduction}

Large Language Models (LLMs) are increasingly used to process unstructured data due to their strong capabilities in semantic understanding and reasoning.
Modern data platforms have begun to develop LLM-powered semantic operators or User-Defined Functions (UDFs), including commercial systems such as Google BigQuery~\cite{google2026bigqueryai}, Snowflake Cortex AI~\cite{snowflake2026cortexai}, and Alibaba PolarDB for AI~\cite{alibaba2026polardbai}, as well as prototype systems for semantic query processing~\cite{DBLP:journals/pvldb/PatelJPGAGZ25}. This integration poses new problems for physical optimization in database systems. 
Unlike classical relational operators, which are CPU- or I/O-centric, LLM-powered semantic operators require LLM invocations over individual records or batches of records. 
Each invocation incurs substantial latency and monetary token cost. Consequently, the primary optimization objective shifts from minimizing CPU operators or I/O transfer to reducing LLM invocations.  This shift calls for new physical algorithms designed specifically for LLM-based semantic operators.

Among semantic operators, semantic \textsf{ORDER BY} supports queries whose ordering criterion is expressed in natural language.
For example, ordering reviews by their helpfulness for a purchase decision, passages by the strength of their evidential support, or retrieved documents by their relevance to a user question. 
Such functionality is central to document reranking in retrieval-augmented generation (RAG)~\cite{DBLP:conf/nips/YuP0WYZSC24} and semantic search. As semantic \textsf{ORDER BY} requires an LLM to infer a global ordering over input items, when the input exceeds the effective context window of the LLM, no single invocation can inspect all items and produce a global ordering. 
A natural solution is divide-and-conquer sorting: inferring the local order of item subsets that fit in a context window, and combining local orders into a global order. 
This makes processing Semantic \textsf{ORDER BY} substantially more difficult than processing record-independent operators such as semantic filtering. 


Existing LLM-based ranking methods can be broadly categorized as pointwise, pairwise, and listwise according to their prompting strategies. 
Pointwise methods ask the LLM to score each item independently and then sort the generated scores~\cite{DBLP:conf/emnlp/DrozdovZD0RWAIM23, DBLP:conf/emnlp/SachanLJAYPZ22}. 
Although easy to parallelize, these methods suffer from poor score calibration across independent calls. 
Pairwise methods invoke the LLM to compare two items at a time and aggregate the results through sorting or graph-based algorithms~\cite{DBLP:conf/acl/LuoCH024, DBLP:conf/naacl/QinJHZWYSLLMWB24}. 
These methods directly model pairwise relationships of items but underutilize the context window and can require prohibitively many invocations. 
%
Listwise methods place multiple items into one prompt and request a ranked permutation~\cite{DBLP:conf/sigir/ZhuangZKZ24, DBLP:conf/emnlp/0001YMWRCYR23, DBLP:journals/corr/abs-2603-17223, DBLP:journals/corr/abs-2305-02156}. By amortizing one invocation over multiple item comparisons, listwise prompting better utilizes the model’s context window. 
This balance between ranking quality and invocation efficiency makes listwise prompting serve as the key primitive in semantic \textsf{ORDER BY} algorithms~\cite{DBLP:journals/corr/abs-2603-17223, DBLP:journals/corr/abs-2509-00303}. 
Under a fixed prompt template, effective context length, and an expected item serialization length, the number of listwise invocations serves as a near-linear proxy for token cost. 

Recent work has explored two main ways to organize listwise calls for semantic ranking. Pivot-based algorithms recursively route items into buckets around selected pivots~\cite{DBLP:journals/corr/abs-2603-17223}, while merge-based algorithms first construct locally sorted runs and then integrate them~\cite{DBLP:journals/corr/abs-2509-00303}. From the perspective of listwise utilization, a merge-based design aligns better than a pivot-based design, particularly in full ranking. Pivot-based partitioning primarily uses a listwise result to assign items to buckets relative to pivots, while the relationships among non-pivots still need to be resolved in subsequent recursions. In contrast, the merge operation can preserve all partial orders from listwise calls in sorted runs during merge passes.

Despite this advantage, the merge phase of listwise semantic \textsf{ORDER BY} has received limited attention. 
Prior work proposes an initial two-way external merge algorithm~\cite{DBLP:journals/corr/abs-2509-00303}. Given $n$ items, the algorithm first generates $r = \lceil n/m \rceil$ sorted runs and then combines them through  $ \lceil \log_2 r \rceil$ merge passes, where each pass processes the current runs many times listwise. 
A natural improvement is to generalize this two-way merge to $k$-way merge, reducing the number of merge passes to $ \lceil \log_k r \rceil$. However, due to the interplay with listwise calls, this generalization is much more complicated than that for classical external merge sort. 
An essential problem arises: \emph{how to effectively exploit listwise calls to merge multiple sorted runs?} 
Addressing this problem entails suitable data structures to manage the intermediate results through the interaction with listwise calls, and more importantly, output a safe and consistent ordering. 
Furthermore, unlike traditional $k$-way merge, there is an underlying paradox in the merge fan-in $k$. For semantic \textsf{ORDER BY}, enlarging $k$ reduces the number of merge passes; however, it will increase the invocations of LLM within each merge pass, under a given listwise batch size $m$, because one group will process more sorted runs.
In practice, simply pursuing a large fan-in $k$ may reduce the expected item throughput of listwise calls. As a result, a principled method to select an optimal $k$ is necessary for achieving effective utilization of tokens of LLMs. 

To solve these problems, this paper proposes cost-effective multi-way merge algorithms for semantic \textsf{ORDER BY} under prompting listwise LLM ranking. First, we formally adapt classical $k$-way merge framework to the processing of semantic \textsf{ORDER BY}. In this framework, two algorithms respectively adapt the classical sliding window merge and loser tree merge to the interaction of listwise calls, managing intermediate sorting results from $k$ sorted runs. 
Specifically, the sliding window merge packs the heads of sorted runs into one listwise context and emits only the `safe prefix' in the result.
In contrast, the loser tree merge explicitly maintains global orderedness across the representatives of $k$ runs, where one item emits calls for tree path refresh with one listwise invocation, enabling more robust ranking over skew item distribution. 
Second, we unify loser-tree merge as LT($h$), where segment length $h$ controls how many items one refresh emits ($h{=}1$ is the classical case). We also extend these algorithms to \LIMIT~$K$ query through pipeline-parallelized execution, avoiding unnecessary materialization of intermediate ranking. 
Finally, to determine an optimal merge fan-in $k$, we develop algorithm-specific cost models that estimate the number of listwise calls with regard to $k$. These models allow users to select the fan-in 
$k$ that minimizes total LLM cost in merge passes, without side effects on ranking quality.
We conduct substantial experiments on 6 benchmarks including 155 queries, comparing against 7 baseline algorithms among pointwise, pairwise, and listwise methods. Measured gains over listwise baselines are reported only after the revised experiments complete; layout-preview figures in this checkpoint are not experimental results. 
Overall, our algorithms and cost models provide a principled algorithmic solution in semantic \textsf{ORDER BY} for modern AI-oriented data platforms.